% GENERATED FILE. Edit canonical lessons, then run npm run generate:books.
% canonical-source-hash: fnv1a64:1e6e4069180a4797
% canonical-lessons: RU-C512-marafon, RU-C512-begun, RU-C512-finish, RU-C512-otstavat, RU-C512-distantsiya

\chapter{The Marathon}
\label{ch:ru-the-marathon}
\hlchaptermodality{512}

\begin{chapteropening}
\textbf{By the end of this chapter:} \emph{I can talk about running a race.}

Complete the last lesson of chapter 512: I can talk about running a race.
\end{chapteropening}

\section[\texorpdfstring{\ru{марафон} (marafón)}{марафон (marafón)}]{\ru{марафон} (marafón) — a marathon}

\label{lesson:RU-C512-marafon}



\begin{quote}
Before the new one: say the Russian for a medal, then the Russian for to award.
\end{quote}

\subsection*{You'll want to know: \ru{марафон}}
\textbf{\ru{марафон}} (\emph{marafón}) — \textquotedblleft{}a marathon\textquotedblright{}.

\begin{quote}
\ru{Мой} \ru{брат} \ru{пробежал} \ru{свой} \ru{первый} \ru{марафон}.
\end{quote}

\begin{quote}
\emph{Moy brat probezhál svoy pérvyy marafón.}
\end{quote}

\begin{quote}
My brother has run his first marathon.
\end{quote}

Masculine.

\begin{cousinweb}
\emph{\ru{марафон}} was borrowed from Greek \emph{Marath\'{\={o}}n} (EtymDB), a place name built on \emph{márathon} 'fennel'.
\end{cousinweb}

\begin{grammarlens}[title={What you've built}]
One more word for telling, explaining and describing at B1.
\end{grammarlens}

\subsection*{Your turn}
\begin{itemize}
\raggedright
  \item \emph{Say it:} \emph{marafón}
  \item \emph{Say it:} \emph{marafón}, once more
  \item \emph{Say it:} a sentence of your own with \emph{marafón}
  \item \emph{Recall:} say the Russian for a medal, then the Russian for to award, then say \emph{marafón} again
\end{itemize}

\begin{tcolorbox}[breakable,colback=teal!4,colframe=teal!35!black,title={Before you move on}]
What does \ru{марафон} mean? (\textquotedblleft{}A marathon\textquotedblright{}.) Say it once more.
\end{tcolorbox}
\section[\texorpdfstring{\ru{бегун} (begún)}{бегун (begún)}]{\ru{бегун} (begún) — a runner}

\label{lesson:RU-C512-begun}



\begin{quote}
Before the new one: say the Russian for to award, then the Russian for a marathon.
\end{quote}

\subsection*{You'll want to know: \ru{бегун}}
\textbf{\ru{бегун}} (\emph{begún}) — \textquotedblleft{}a runner\textquotedblright{}.

\begin{quote}
\ru{Бегуны} \ru{ждут} \ru{сигнала} \ru{на} \ru{старте}.
\end{quote}

\begin{quote}
\emph{Beguný zhdut signála na stárte.}
\end{quote}

\begin{quote}
The runners are waiting for the signal at the start.
\end{quote}

Masculine; genitive \ru{бегуна}.

\begin{cousinweb}
\emph{\ru{бегун}} is built on \emph{\ru{бег}} 'running' with \emph{-\ru{ун}}; EtymDB has no entry.
\end{cousinweb}

\begin{grammarlens}[title={What you've built}]
One more word for telling, explaining and describing at B1.
\end{grammarlens}

\subsection*{Your turn}
\begin{itemize}
\raggedright
  \item \emph{Say it:} \emph{begún}
  \item \emph{Say it:} \emph{begún}, once more
  \item \emph{Say it:} a sentence of your own with \emph{begún}
  \item \emph{Recall:} say the Russian for to award, then the Russian for a marathon, then say \emph{begún} again
\end{itemize}

\begin{tcolorbox}[breakable,colback=teal!4,colframe=teal!35!black,title={Before you move on}]
What does \ru{бегун} mean? (\textquotedblleft{}A runner\textquotedblright{}.) Say it once more.
\end{tcolorbox}
\section[\texorpdfstring{\ru{финиш} (fínish)}{финиш (fínish)}]{\ru{финиш} (fínish) — the finish}

\label{lesson:RU-C512-finish}



\begin{quote}
Before the new one: say the Russian for a marathon, then the Russian for a runner.
\end{quote}

\subsection*{You'll want to know: \ru{финиш}}
\textbf{\ru{финиш}} (\emph{fínish}) — \textquotedblleft{}the finish\textquotedblright{}.

\begin{quote}
\ru{Перед} \ru{финишем} \ru{он} \ru{обогнал} \ru{всех}.
\end{quote}

\begin{quote}
\emph{Péred fínishem on obognál vsekh.}
\end{quote}

\begin{quote}
Before the finish he overtook everyone.
\end{quote}

Masculine.

\begin{cousinweb}
\emph{\ru{финиш}} was borrowed from English \emph{finish} (EtymDB).
\end{cousinweb}

\begin{grammarlens}[title={What you've built}]
One more word for telling, explaining and describing at B1.
\end{grammarlens}

\subsection*{Your turn}
\begin{itemize}
\raggedright
  \item \emph{Say it:} \emph{fínish}
  \item \emph{Say it:} \emph{fínish}, once more
  \item \emph{Say it:} a sentence of your own with \emph{fínish}
  \item \emph{Recall:} say the Russian for a marathon, then the Russian for a runner, then say \emph{fínish} again
\end{itemize}

\begin{tcolorbox}[breakable,colback=teal!4,colframe=teal!35!black,title={Before you move on}]
What does \ru{финиш} mean? (\textquotedblleft{}The finish\textquotedblright{}.) Say it once more.
\end{tcolorbox}
\section[\texorpdfstring{\ru{отставать} (otstavát)}{отставать (otstavát)}]{\ru{отставать} (otstavát) — to fall behind}

\label{lesson:RU-C512-otstavat}



\begin{quote}
Before the new one: say the Russian for a runner, then the Russian for the finish.
\end{quote}

\subsection*{You'll want to know: \ru{отставать}}
\textbf{\ru{отставать}} (\emph{otstavát}) — \textquotedblleft{}to fall behind\textquotedblright{}.

\begin{quote}
\ru{Я} \ru{часто} \ru{отстаю} \ru{от} \ru{группы}.
\end{quote}

\begin{quote}
\emph{Ya chásto otstayú ot grúppy.}
\end{quote}

\begin{quote}
I often fall behind the group.
\end{quote}

Imperfective; present \ru{отстаю}, \ru{отстаёшь} (-\ru{ава}- becomes -\ru{а}-). Perfective \ru{отстать}. Of a clock, to be slow.

\begin{cousinweb}
\emph{\ru{отставать}} is built on \emph{\ru{от}-} 'away' and \emph{\ru{ставать}}, the stem of \emph{\ru{стать}} 'to stand, to become'; EtymDB lists the verb without a source.
\end{cousinweb}

\begin{grammarlens}[title={What you've built}]
One more word for telling, explaining and describing at B1.
\end{grammarlens}

\subsection*{Your turn}
\begin{itemize}
\raggedright
  \item \emph{Say it:} \emph{otstavát}
  \item \emph{Say it:} \emph{otstavát}, once more
  \item \emph{Say it:} a sentence of your own with \emph{otstavát}
  \item \emph{Recall:} say the Russian for a runner, then the Russian for the finish, then say \emph{otstavát} again
\end{itemize}

\begin{tcolorbox}[breakable,colback=teal!4,colframe=teal!35!black,title={Before you move on}]
What does \ru{отставать} mean? (\textquotedblleft{}To fall behind\textquotedblright{}.) Say it once more.
\end{tcolorbox}
\section[\texorpdfstring{\ru{дистанция} (distántsiya)}{дистанция (distántsiya)}]{\ru{дистанция} (distántsiya) — a distance (in a race)}

\label{lesson:RU-C512-distantsiya}



\begin{quote}
Before the new one: say the Russian for the finish, then the Russian for to fall behind.
\end{quote}

\subsection*{You'll want to know: \ru{дистанция}}
\textbf{\ru{дистанция}} (\emph{distántsiya}) — \textquotedblleft{}a distance (in a race)\textquotedblright{}.

\begin{quote}
\ru{Она} \ru{бегает} \ru{на} \ru{длинные} \ru{дистанции}.
\end{quote}

\begin{quote}
\emph{Oná bégayet na dlínnyye distántsii.}
\end{quote}

\begin{quote}
She runs long-distance races.
\end{quote}

Feminine.

\begin{cousinweb}
\emph{\ru{дистанция}} was borrowed from Latin \emph{distantia} 'distance' (EtymDB), from \emph{distāre} 'to stand apart' (Lewis \& Short).
\end{cousinweb}

\begin{grammarlens}[title={What you've built}]
One more word for telling, explaining and describing at B1.
\end{grammarlens}

\subsection*{Your turn}
\begin{itemize}
\raggedright
  \item \emph{Say it:} \emph{distántsiya}
  \item \emph{Say it:} \emph{distántsiya}, once more
  \item \emph{Say it:} a sentence of your own with \emph{distántsiya}
  \item \emph{Recall:} say the Russian for the finish, then the Russian for to fall behind, then say \emph{distántsiya} again
\end{itemize}

\begin{tcolorbox}[breakable,colback=teal!4,colframe=teal!35!black,title={Before you move on}]
What does \ru{дистанция} mean? (\textquotedblleft{}A distance (in a race)\textquotedblright{}.) Say it once more.
\end{tcolorbox}
